\section{Introduction}
\label{sec:intro}

Recommenders built on large language models (LLMs) routinely reduce recommendation to a binary
judgment---``Would this user like this item? Answer Yes or No''---and read the probability of \emph{Yes} as the
recommendation score~\citep{bao2023tallrec,zhang2025collm}. The same number is then implicitly treated as the
model's \emph{confidence}: a high \py{} is taken to mean that the model is sure the user will like the item.
Whether that reading is justified has, to our knowledge, never been tested against labels that ask the same
question. Prior work quantifies the uncertainty of list-wise LLM recommenders~\citep{kweon2025uqrec}, shapes
training rewards with confidence~\citep{fan2026ugr}, or audits verbalized confidence of zero-shot models against
catalog membership~\citep{ravikumar2026hallucinating}; none asks \emph{what the yes/no confidence of a
personalised LLM recommender is made of}.

We study six questions that practitioners ask of such a system:
(Q1) when the model is right, is it sure? (Q2) when it is wrong, is that because it was unsure---or is it
confidently wrong? (Q3) does acting on high confidence over-expose popular items? (Q4) how well does the
yes/no confidence track ground truth? (Q5) are popular items systematically more confidently endorsed than niche
ones, and is that endorsement \emph{warranted}? (Q6) can confidence identify noisy training data?

\paragraph{Key observation.} \DATANEEDED{Pilot-1/B3 result: a(i) spread, popularity link, share of confident
errors explained by the item term; pseudonym knockout}. Our central hypothesis is that \py{} decomposes into an
\emph{item term}---a willingness to say \emph{Yes} to an item regardless of who asks, which tracks how familiar
or popular the item is---and a \emph{personal-evidence term}. Only the latter is a recommendation signal.
Crucially, the item term is not removed by within-user calibration: any monotone recalibration of \py{} leaves a
user's ranking unchanged~\citep{kweon2022calibrated}, whereas an item-dependent offset reorders it.

\paragraph{Method.} \method{} asks the mirrored question (``Would this user \emph{dislike} this item?'') and
scores the difference of the two log-odds. The item term enters both answers with the same sign and cancels;
user-specific evidence enters with opposite signs and is doubled. The two-prompt symmetrization estimator itself
is not new~\citep{burns2023ccs}; our contribution is to show that in \emph{personalised} recommendation the
artifact it removes is item-specific and popularity-graded, that removing it changes rankings in a way that
equal-compute prompt ensembles and item-prior subtraction do not~\citep{zhao2021calibrate,holtzman2021surface},
and that mirror tuning keeps the benefit after fine-tuning.

\paragraph{Contributions.}
(1) The first calibration study of token-level yes/no confidence of LLM recommenders against rated labels that
match the question, answering Q1--Q5 on \DATANEEDED{4 datasets $\times$ 2 backbones}, with a validity check
against fully observed feedback~\citep{gao2022kuairec}.
(2) Evidence that the confidence contains an item-level, popularity-graded acquiescence term, including a
causal pseudonym-knockout test with a real-brand placebo.
(3) \method{} and mirror tuning, with \DATANEEDED{main gains}.
(4) A characterised set of what does \emph{not} work (within-user recalibration, uncertainty-weighted pruning),
reported with multi-seed controls (Appendix).
